\begin{textsample}{POS Dim 2 – 2018\_c5 – Score -1.00 – 2018\_c5/enade\_2018\_serv\_soc\_f2ecfc474e.txt}  \label{ex:f2_pos_316}
CAPÍTULO 5 DISTRIBUIÇÃO DOS CONCEITOS 5.1 PANORAMA NACIONAL DA DISTRIBUIÇÃO DOS CONCEITOS O Conceito Enade é um indicador de qualidade que avalia os cursos por intermédio dos desempenhos dos estudantes a partir dos resultados do Enade.
Os valores possíveis vão de 1 ( pior situação ) a 5 ( melhor situação ).
Os cursos com apenas um ou sem nenhum concluinte participante não obtêm o Conceito Enade, ficando Sem Conceito ( SC ).
Na Tabela 5.1, são apresentadas a quantidade e a distribuição de cursos de Serviço Social participantes do Enade/2018, por faixa de conceito e Grande Região.
Enfatiza -se, mais uma vez, que a diferença entre os cursos tabulados neste capítulo e no capítulo 2 corresponde aos cursos Sem Conceito, ou seja àqueles, em princípio, sem alunos concluintes que houvessem participado da prova ou àqueles que tivessem tão somente um aluno concluinte19.
Observando -se os dados da Tabela 5.1, nota -se que, dos 330 cursos participantes, 128 ( 38,8 \% ) classificaram -se com conceito 3, o valor modal.
Esse foi também o conceito modal nas regiões : Norte ( 35,5 \% ), Nordeste ( 37,4 \% ) e Sudeste ( 43,4 \% ).
O conceito 2 foi o segundo mais frequente em nível nacional ( 27,0 \%, correspondendo a 89 cursos ), e o conceito 4, o terceiro ( 17,6 \%, correspondendo a 58 cursos ).
Houve, ainda, 22 cursos ( 6,7 \% ) que receberam conceito 5, e outros 22 cursos ( 6,7 \% ) que receberam conceito 1.
Dos 330 cursos de Serviço Social, 11 cursos ( 3,3 \% ) ficaram Sem Conceito ( SC ).
A região Norte participou com 31 cursos ou 9,4 \% do total nacional.
Desses, 12 cursos receberam o conceito 2, o que equivale a 38,7 \% do total regional, além de ser o conceito modal na região.
A 11 cursos ( 35,5 \% ) atribuiu -se conceito 3, e a três cursos ( 9,7 \% ), conceito 4. 19 Esses cursos com somente um concluinte não foram considerados no capítulo 2. 4.
Cinco cursos ( 16,1 \% ) receberam conceito 1 e nenhum curso ficou Sem Conceito ( SC ) ou recebeu conceito 5.
A região Nordeste participou com 107 cursos ou 32,4 \% do total nacional.
Nessa região, 40 cursos ( 37,4 \% em termos regionais ) obtiveram conceito 3, o conceito modal para a região, como já comentado.
Os conceitos 2 e 4 foram atribuídos, respectivamente, a 39 e a 12 cursos ( 36,4 \% e 11,2 \%, respectivamente ).
Ainda nessa região, nove cursos ( 8,4 \% ) foram avaliados com o conceito 5 e seis cursos ( 5,6 \% ), com conceito 1.
Nenhum curso ficou Sem Conceito ( SC ) nessa região.
Dos 129 cursos participantes da região Sudeste ( 39,1 \% do total nacional ), o conceito 3 foi o conceito modal, como já anteriormente assinalado, tendo sido obtido por 56 cursos ( 43,4 \% ).
O conceito 1 foi atribuído a dez cursos ( 7,8 \% ), e o conceito 2, a 25 cursos ( 19,4 \% ).
Receberam os conceitos 4 e 5, respectivamente, 25 cursos ( 19,4 \% ) e cinco cursos ( 3,9 \% ).
Nessa região, oito cursos ficaram Sem Conceito ( SC ).
Os 44 cursos da região Sul corresponderam a 13,3 \% do total nacional.
Houve predominância do conceito 4, atribuído a 36,4 \% dos cursos da região ( conceito modal ), ou seja, atribuído a 16 dos 44 cursos participantes na região Sul.
O conceito 1 foi atribuído a um curso ( 2,3 \% ) e os conceitos 2 e 3, a seis cursos ( 13,6 \% ) e 15 cursos ( 34,1 \% ), respectivamente.
Cinco cursos ( 11,4 \% ) receberam conceito 5 e um curso ficou Sem Conceito ( SC ) na região Sul.
Já dos 19 cursos participantes na região Centro-Oeste ( 5,8 \% do total nacional ), sete ( 36,8 \% em termos regionais ) receberam conceito 2, o conceito modal.
Seis cursos ( 31,6 \% ) obtiveram conceito 3, e outros dois ( 10,5 \% ), conceito 4.
Além disso, três cursos ( 15,8 \% ) receberam conceito 5 e um curso ficou Sem Conceito ( SC ).
Nessa região, a nenhum curso foi atribuído o conceito1.
No Gráfico 5.1, é apresentada a distribuição cumulativa do Conceito Enade segundo Grande Região.
Em um gráfico de distribuição cumulativa, em cada valor de abscissa ( eixo x ), o valor apresentado no eixo das ordenadas ( eixo y ) é o valor acumulado da variável.
Por exemplo, a região Nordeste ( linha verde ) apresenta um pouco mais de 80,0 \% dos cursos com conceito 3 ou abaixo deste valor.
A região Sul ( linha rosa ), por outro lado, apresenta menos de 52,0 \% dos cursos na mesma situação : com conceito 3 ou abaixo.
Na região Norte ( linha azul ), todos os cursos têm conceito 4 ou abaixo disso, já que a poligonal cumulativa alcança 100,0 \% para o conceito 4.
Em linhas gerais, em um gráfico de distribuição cumulativa, poligonais mais à esquerda representam Grandes Regiões com pior distribuição, e poligonais mais à direita, Grandes Regiões com melhores conceitos.
A região Norte ( linha azul ) apresenta os cursos com conceitos mais baixos, ao passo que a região Sul ( linha rosa ) apresenta as distribuições com valores maiores.
Cursos em IES situadas nas regiões Sudeste ( linha vermelha ), Nordeste ( linha verde ) e Centro-Oeste ( linha laranja ) apresentam uma situação intermediária. 5.2 CONCEITOS POR CATEGORIA ADMINISTRATIVA, POR MODALIDADE DE ENSINO E GRANDE REGIÃO Na Tabela 5.2, é apresentada a distribuição dos cursos de Serviço Social participantes do Enade/2018, por Categoria Administrativa e Modalidade de Ensino, de acordo com os conceitos por eles alcançados, segundo as Grandes Regiões brasileiras.
Dos 330 cursos participantes, 67 ( 20,3 \% ) eram ministrados em Instituições Públicas, e 263 ( 79,7 \% ), em Instituições Privadas.
Quanto à Modalidade de Ensino, 24 cursos eram a Distância, e 306, Presenciais.
De acordo com as informações da Tabela 5.2, em termos nacionais, entre Instituições Públicas e Privadas, observa -se que, dos 22 cursos avaliados com conceito 5, 16 eram oferecidos em IES Públicas, e seis, em IES Privadas.
Dos 67 cursos participantes de IES Públicas, o conceito 4 foi o valor modal, atribuído a 30 cursos ( 44,8 \% ).
Dos demais cursos de Instituições Públicas participantes, um curso foi avaliado com conceito 1, cinco cursos ( 7,5 \% ) foram avaliados com conceito 2, 15 cursos, com conceito 3, e 16 cursos, com conceito 5.
Nenhum curso ficou Sem Conceito ( SC ) nesta categoria.
Na rede privada, o conceito modal foi 3, atribuído a 113 cursos, o correspondente a 43,0 \% dos 263 cursos da categoria.
Dos demais cursos participantes, 21 ( 8,0 \% ) receberam conceito 1, e 84 ( 31,9 \% ), conceito 2.
O conceito 4 foi atribuído a 28 cursos ( 10,6 \% ), e o conceito 5, a seis cursos ( 2,3 \% ).
Nesta Categoria Administrativa, 11 cursos ( 4,2 \% ) ficaram Sem Conceito ( SC ).
Dos 306 cursos participantes da Modalidade de Ensino Presencial, o conceito 3 foi o valor modal, atribuído a 124 cursos ( 40,5 \% ).
Dos demais cursos nesta Modalidade de Ensino, 16 cursos ( 5,2 \% ) receberam conceito 1, 77 cursos ( 25,2 \% ) receberam conceito 2, 58 cursos ( 19,0 \% ), conceito 4, e 22 cursos ( 7,2 \% ), conceito 5.
Nesta Modalidade de Ensino, nove cursos ( 2,9 \% ) ficaram Sem Conceito ( SC ).
Aos 24 cursos a Distância informados atribuíram -se os conceitos 1 ( seis cursos ), 2 ( 12 cursos, valor modal ) e 3 ( quatro cursos ).
Nessa Modalidade de Ensino, dois cursos ficaram Sem Conceito ( SC ) e nenhum curso recebeu os conceitos 4 ou 5.
Na análise por região, observa -se que, na região Norte, as Instituições Privadas participaram com 25 cursos ( 80,6 \% do total regional ), dos quais nenhum ficou Sem Conceito ( SC ) ou recebeu conceito 5.
O conceito modal para as Instituições Privadas na região foi 2, atribuído a 11 cursos, correspondendo a 44,0 \% dessa categoria na região.
Cinco cursos foram avaliados com conceito 1, oito cursos, com conceito 3, e ainda um curso, com conceito 4.
As Instituições Públicas participaram com seis cursos na região Norte ( 19,4 \% do total regional ), aos quais foram atribuídos os conceitos 2, para um curso, conceito 3, para três cursos, e conceito 4, para dois cursos.
Nenhum curso ficou Sem Conceito ( SC ) ou recebeu os conceitos 1 ou 5 nessa combinação de categoria e região.
Dos 31 cursos oferecidos na região Norte, 30 eram presenciais.
Para esses, o conceito modal 2 recebeu 12 cursos e nenhum curso foi atribuído ao conceito 5 ou ficou Sem Conceito ( SC ).
Para os demais cursos desta Modalidade de Ensino foram alocados os conceitos 1 ( quatro cursos ), 3 ( 11 cursos ) e 4 ( três cursos ).
O único curso a Distância informado, foi avaliado com conceito 1.
Na região Nordeste, a rede privada concentrou 89 dos 107 cursos participantes, o equivalente a 83,2 \% do total da região.
Dentre os cursos oferecidos por IES Privadas no Nordeste, 39 cursos foram avaliados com conceito 2, o conceito modal. 37 cursos obtiveram conceito 3, cinco cursos receberam conceito 4, um curso, conceito 5, e ainda seis cursos, conceito 1.
Nessa combinação de categoria e região, um ficou Sem Conceito ( SC ).
As Instituições Públicas da região Nordeste participaram com 18 cursos ( 16,8 \% do total da região ).
Desses, oito obtiveram conceito 5, conceito modal, seguidos por outros sete cursos, que receberam conceito 4.
Três cursos receberam conceito 3, e nenhum curso ficou Sem Conceito ( SC ) ou recebeu os demais conceitos nessa região.
Dos cursos oferecidos na região Nordeste 103 eram de Educação Presencial, tendo o conceito 3 como modal, com 39 cursos atribuídos.
Os demais cursos foram alocados aos conceitos 1 ( seis cursos ), 2 ( 36 cursos ), 4 ( 12 cursos ) e 5 ( nove cursos ).
Um curso ficou Sem Conceito ( SC ).
Dos quatro cursos a Distância informados, o conceito modal 2 foi atribuído a três cursos, e um curso foi alocado ao conceito 3.
Nesta Modalidade de Ensino nenhum curso ficou Sem Conceito ( SC ) ou recebeu os conceitos 1, 4 ou 5.
Na região Sudeste, a proporção de cursos da rede privada, 80,6 \%, correspondeu a 104 dos 129 cursos participantes.
Nessa categoria e região, o conceito modal foi 3, atribuído a 51 cursos, e oito cursos ficaram Sem Conceito ( SC ).
Os demais cursos foram avaliados com conceito 1 ( nove cursos ), conceito 2 ( 23 cursos ), conceito 4 ( 12 cursos ) e conceito 5 ( um curso ).
Dos 25 cursos oferecidos por Instituições Públicas na região Sudeste, o conceito modal foi 4, atribuído a 13 cursos.
Nessa região, a um curso foi atribuído o conceito 1, a dois cursos, o conceito 2, a cinco cursos, o conceito 3, e ainda a quatro cursos, o conceito 5.
Nenhum curso de IES Públicas ficou Sem Conceito ( SC ).
Também na região Sudeste, a maioria dos cursos foi oferecido na modalidade Presencial : 118 cursos.
Para esses, a moda foi o conceito 3, com 55 cursos.
Os demais cursos presenciais da região foram atribuídos pelos conceitos 1 ( cinco cursos ), 2 ( 21 cursos ), 4 ( 25 cursos ) e 5 ( cinco cursos ).
Ficaram Sem Conceito ( SC ) sete cursos.
Na região Sudeste, os dez cursos oferecidos na modalidade Educação a Distância obtiveram os conceitos 1 ( cinco cursos ), 2 ( quatro cursos ) e 3 ( um curso ).
Nesta Modalidade de Ensino um curso ficou Sem Conceito ( SC ) e nenhum curso recebeu os conceitos 4 ou 5.
As Instituições Privadas concentraram 30 dos 44 cursos participantes da região Sul, 68,2 \% do total regional.
Desses, 11 cursos obtiveram conceito 3, o conceito modal.
Nesta combinação de Categoria Administrativa e Grande Região, um curso recebeu conceito 1, outros cinco, conceito 2, nove cursos, conceito 4, e ainda outros três cursos, conceito 5.
Um curso ficou Sem Conceito ( SC ).
As Instituições Públicas na região Sul participaram com 14 cursos ( 31,8 \% dos cursos da região ), aos quais foram atribuídos os conceitos 2 ( um curso ), 3 ( quatro cursos ), 4 ( sete cursos, conceito modal ) e 5 ( dois cursos ).
Nenhum curso da região Sul oferecido por IES Públicas ficou Sem Conceito ( SC ) ou recebeu o conceito 1.
Dos cursos oferecidos na região Sul, 38 eram de Educação Presencial, e seis, de Educação a Distância.
O conceito modal para os cursos presenciais foi 4, com 16 cursos.
Os demais cursos receberam os conceitos 1 ( um curso ), 2 ( três cursos ), 3 ( 13 cursos ) e 5 ( cinco cursos ).
Nessa Modalidade de Ensino nenhum curso ficou Sem Conceito ( SC ).
Dos seis cursos a distância, um curso ficou Sem Conceito ( SC ) e a nenhum curso foram atribuídos os conceitos 1, 4 e 5.
Foi atribuído o conceito 2, conceito modal, a três cursos e o conceito 3, a dois cursos.
Na região Centro-Oeste, 15 dos 19 cursos participantes eram de Instituições Privadas ( 78,9 \% em termos regionais ).
Desses, seis receberam conceito 2, e outros seis, conceito 3.
Os conceitos 4 e 5 receberam um curso, cada, e nenhum curso recebeu conceito 1.
Nesta categoria e região um curso ficou Sem Conceito ( SC ).
Quanto aos quatro cursos oferecidos por Instituições Públicas na região Centro-Oeste ( 21,1 \% do total regional ), os conceitos 2 e 4 foram atribuídos a um curso, cada, e o conceito 5, conceito modal, atribuídos a dois cursos.
Nenhum curso ficou Sem Conceito ( SC ) ou recebeu os demais conceitos nessa combinação de categoria e região.
Dos 17 cursos oferecidos, na região, na Modalidade Presencial, a seis cursos foi atribuído o conceito modal 3, outros cinco cursos receberam o conceito 2, dois cursos, o conceito 4, e outros três cursos, conceito 5.
Um curso ficou Sem Conceito ( SC ) e a nenhum curso foi atribuído o conceito 1.
Os dois cursos oferecidos na Modalidade a Distância foram avaliados com conceito 2.
No Gráfico 5.2, apresenta -se a distribuição do Conceito Enade, segundo a Categoria Administrativa da IES.
Os cursos em IES Públicas ( linha azul ) apresentam uma poligonal mais à direita do que os em IES Privadas ( linha verde ) e, consequentemente, uma distribuição de Conceitos Enade com valores maiores. 5.3 CONCEITOS POR ORGANIZAÇÃO ACADÊMICA E GRANDE REGIÃO Na Tabela 5.3, encontra -se a distribuição dos conceitos atribuídos aos cursos participantes do Enade/2018, na Área de Serviço Social, por Organização Acadêmica, segundo as Grandes Regiões brasileiras.
Dos 330 cursos de Serviço Social participantes, 146 eram oferecidos em Universidades, 68 em Centros Universitários, 115 em Faculdades e um em CEFET/IFET.
Essa distribuição corresponde a, respectivamente, 44,2 \%, 20,6 \%, 34,8 \% e 0,3 \% dos cursos.
De acordo com os dados apresentados, de todos os 22 cursos avaliados com conceito 5, 17 eram vinculados a Universidades.
Os cursos vinculados a esse tipo de Organização Acadêmica tiveram o conceito 3 como conceito modal, atribuído a 45 cursos ( 31,5 \% ).
Os demais cursos vinculados a Universidades receberam os conceitos 1 ( oito cursos ), 2 ( 29 cursos ), conceito 2, 4 ( 43 cursos ) e conceito 5 ( 17 cursos, como já mencionado ).
Ficaram Sem Conceito ( SC ), quatro cursos.
Dos cursos em Centros Universitários, o conceito modal também foi 3, atribuído a 38 cursos ( 55,9 \% ).
Quatro cursos vinculados a esse tipo de Organização Acadêmica ficaram Sem Conceito ( SC ), e os outros cursos receberam os conceitos 1 ( quatro cursos ), 2 ( 15 cursos ), 4 ( seis cursos ) e 5 ( um curso ).
Dos 115 cursos mantidos por Faculdades, 45 cursos ( 39,1 \% ) receberam conceito 2, e outros 45 cursos ( 39,1 \% ), conceito 3, conceitos modais.
Aos demais cursos foram atribuídos os conceitos 1 ( dez cursos ), 4 ( nove cursos ) e 5 ( três cursos ).
Dos cursos oferecidos em Faculdades, três cursos ficaram Sem Conceito ( SC ).
O único curso mantido por CEFET/IFET, 0,3 \% do total, recebeu conceito 5.
Considerando -se separadamente as regiões brasileiras, verifica -se que, na região Norte, as Universidades participaram com apenas nove dos 31 cursos avaliados, ao qual foram atribuídos os conceitos 2 ( três cursos ), 3 ( quatro cursos, valor modal ) e 4 ( dois cursos ).
Nenhum curso ficou Sem Conceito ( SC ) ou recebeu os demais conceitos nessa combinação de categoria e região.
Os Centros Universitários da região Norte foram representados por sete cursos, os quais obtiveram os conceitos 1 ( dois cursos ), 2 ( quatro cursos ) e 3 ( um curso ).
Nenhum curso ficou Sem Conceito ( SC ) ou recebeu os demais conceitos nessa combinação de categoria e região.
As Faculdades participaram com 15 cursos na região Norte, os quais receberam os conceitos 1 ( três cursos ), 2 ( cinco cursos ), 3 ( seis cursos, conceito modal ) e 4 ( um curso ).
Nenhum curso ficou Sem Conceito ( SC ) ou recebeu conceito 5.
Já os CEFET/IFET não participaram com cursos na região Norte.
Na região Nordeste, as Universidades participaram com 28 dos 107 cursos da Área de Serviço Social oferecidos.
Nenhum curso ficou Sem Conceito ( SC ).
O conceito modal 2 foi atribuído a oito cursos.
Os demais cursos obtiveram os conceitos 1 ( dois cursos ), 3 ( quatro cursos ) e os conceitos 4 e 5 ( sete cursos, cada ).
Os Centros Universitários contaram com 18 cursos participantes na região Nordeste, 12 dos quais receberam o conceito modal 3.
Para os demais cursos, foram atribuídos os conceitos 1 ( um curso ), 2 ( quatro cursos ) e 4 ( um curso ).
Nenhum curso oferecido por Centros Universitários no Nordeste ficou Sem Conceito ( SC ) ou recebeu conceito 5.
As Faculdades foram representadas por 60 cursos na região Nordeste, 27 dos quais receberam conceito 2, o conceito modal.
Os demais cursos receberam os conceitos 1 ( três cursos ), 3 ( 24 cursos ), 4 ( quatro cursos ) e 5 ( um curso ).
Na região Nordeste, um curso mantido por Faculdades ficou Sem Conceito ( SC ).
Os CEFET/IFET participaram com um curso na região Nordeste, avaliado com conceito 5.
Na região Sudeste, as Universidades concentraram 68 dos 129 cursos de Serviço Social da região.
Dos cursos oferecidos em Universidades, nessa região, o conceito modal foi 3, atribuído a 24 cursos.
Os demais cursos receberam os conceitos 1 ( seis cursos ), 2 ( 11 cursos ), 4 ( 19 cursos ) e 5 ( quatro cursos ).
Quatro cursos ficaram Sem Conceito ( SC ).
Os Centros Universitários participaram com 32 cursos na região Sudeste, dos quais 19 obtiveram o conceito modal 3.
O conceito 1 foi atribuído a um curso, o conceito 2, a cinco cursos, o conceito 4, a três cursos, e o conceito 5, a um curso.
Três cursos ficaram Sem Conceito ( SC ) nessa combinação de Organização Acadêmica e região.
Já as Faculdades foram representadas por 29 cursos na região Sudeste, dos quais um ficou Sem Conceito ( SC ) e nenhum curso recebeu conceito 5.
O conceito modal 3 foi atribuído a 13 cursos.
Aos demais cursos foram atribuídos os conceitos 1 ( três cursos ), 2 ( nove cursos ) e 4 ( três cursos ).
Os CEFET/IFET não participaram com cursos nesta região.
Dos 44 cursos da região Sul, 30 eram mantidos por Universidades, para os quais o conceito modal foi 4, atribuído a 13 cursos.
Os demais cursos receberam os conceitos 2 ( quatro cursos ), 3 ( nove cursos ) e 5 ( quatro cursos ).
Nesse tipo de organização, nenhum curso da região Sul ficou Sem Conceito ( SC ) ou recebeu conceito 1.
Dos oito cursos participantes de Centros Universitários da região Sul, apenas um ficou Sem Conceito ( SC ), e nenhum recebeu os conceitos 1 ou 5.
Os cursos restantes obtiveram os conceitos 2 ( um curso ), 3 ( quatro cursos, conceito modal ) e 4 ( dois cursos ).
Quanto aos seis cursos vinculados a Faculdades na região Sul, dois receberam o conceito modal 3.
E os demais conceitos 1, 2, 4 e 5 foram atribuídos a um curso, cada.
Nenhum curso ficou Sem Conceito ( SC ).
Já os CEFET/IFET não participaram com cursos na região Norte.
Na região Centro-Oeste, 11 dos 19 cursos oferecidos eram mantidos por Universidades.
Para esse tipo de Organização Acadêmica, nenhum curso ficou Sem Conceito ( SC ) ou recebeu conceito 1.
O conceito modal 3 foram atribuídos a quatro cursos, o conceito 2, a três cursos, e os conceitos 4 e 5 a dois cursos, cada.
Os Centros Universitários da região Centro-Oeste contaram com três cursos, dois dos quais foram avaliados com conceito 3, conceito modal, e um curso recebeu conceito 2.
Dos cinco cursos oferecidos em Faculdades na região Centro-Oeste, três receberam o conceito modal 2, e um curso obteve conceito 5.
Nessa região, um curso oferecido em Faculdades ficou Sem Conceito ( SC ), e nenhum obteve conceito 1, 3 ou 4.
Já os CEFET/IFET não participaram com cursos na região Centro-Oeste.
No Gráfico 5.3, apresenta -se a distribuição do Conceito Enade segundo a Organização Acadêmica da IES.
O único curso oferecido em CEFET/IFET ( linha laranja ) recebeu conceito 5 e corresponde a uma poligonal mais à direita do que as demais, denotando uma melhor distribuição de Conceitos Enade.
No outro extremo, encontram -se os cursos mantidos por Centros Universitários ( linha verde ) e Faculdades ( linha vermelha ), apresentando as respectivas poligonais mais à esquerda e, consequentemente, distribuição de Conceitos Enade com menores valores.
Cursos em Universidades ( linha azul ) apresentam uma situação intermediária.

% matched lemmas: 
\end{textsample}
